% BHU PYQ -- Manually transcribed & error-corrected
% Paper: MTB-102 : Geometry
% Exam: B.Sc. (Hons.) Semester I Examination, 2016-17

\documentclass[11pt,a4paper]{article}
\usepackage[a4paper,top=2.5cm,bottom=2.5cm,left=2.8cm,right=2.8cm,headheight=14pt]{geometry}
\usepackage{amsmath,amssymb}
\usepackage[T1]{fontenc}
\usepackage[utf8]{inputenc}
\usepackage{lmodern}
\usepackage{microtype}
\usepackage{fancyhdr}
\usepackage{enumitem}
\usepackage{hyperref}
\usepackage{lastpage}
\usepackage{parskip}

\hypersetup{colorlinks=true,urlcolor=black,linkcolor=black,
  pdftitle={MTB-102: Geometry -- BHU B.Sc. Sem I 2016-17}}
\pagestyle{fancy}\fancyhf{}
\renewcommand{\headrulewidth}{0.4pt}\renewcommand{\footrulewidth}{0.4pt}
\fancyhead[L]{\small   \textit{Banaras Hindu University}}
\fancyhead[R]{\small   \textit{MTB-102: Geometry}}
\fancyfoot[C]{\small Page  \thepage\ of \pageref{LastPage}}


\newlist{parts}{enumerate}{2}
\setlist[parts,1]{label=(\alph*),leftmargin=*,itemsep=5pt,topsep=3pt}
\setlist[parts,2]{label=(\roman*),leftmargin=*,itemsep=3pt,topsep=2pt}

\newcommand{\pts}[1]{\hfill   \textit{[#1]}}


\begin{document}

\begin{center}
  {\large\bfseries BANARAS HINDU UNIVERSITY}\\[2pt]
  {
\normalsize B.Sc. (Hons.) --- Semester I Examination, 2016-17}\\[6pt]
  \rule{0.8\linewidth}{0.5pt}\\[6pt]
  {\large\bfseries Paper No. MTB-102: Geometry}\\[4pt]
  {
\normalsize Time: 3 Hours \hfill Full Marks: 70}
\end{center}
\rule{\linewidth}{0.4pt}
\smallskip



\noindent   \textit{Instructions: Answer   \textbf{five} questions including Question~1 which is compulsory. Figures in right-hand margin indicate marks.}
\bigskip




\noindent   \textbf{Question 1.}   \textit{(Compulsory --- 2 marks each.)}

\begin{parts}
  \item Find the latus rectum and the point on the conic $\frac{l}{r} = 3 - 8\cos \theta$ whose radius vector is 2.
  \item Find the length of the focal chord of a parabola which passes through a point whose vectorial angle is $\alpha$.
  \item Find the intercepts of the plane $2x - 3y + z = 12$ on the coordinate axes.
  \item Find the coordinates of the foot of the perpendicular from the point $(1, 3, 4)$ upon the plane $2x - y + 2z = 0$.
  \item Write the equation of straight line $x - y = 0, z = 1$ in symmetrical form.
  \item Obtain the condition for two spheres to be orthogonal.
  \item Find the equation of the sphere through the circle $x^2 + y^2 + z^2 = 9$, $2x + 3y + 4z = 5$ and the point $(1, 2, 3)$.
  \item If the vertex of a right circular cone be at the origin and the axis be the $z$-axis, then find the equation of the cone with semi-vertical angle $ \theta$.
  \item Give the definition of reciprocal cones.
  \item Define a right circular cylinder.
  \item Show that five normals can be drawn to a paraboloid from a given point.
\end{parts}

\medskip\rule{\linewidth}{0.2pt}

\medskip


\noindent   \textbf{Question 2.}

\begin{parts}
  \item Find the polar equation of the normal to a conic $\frac{l}{r} = 1 - e\cos \theta$ at a point whose vectorial angle is $\alpha$. \pts{6}
  \item If the tangents at any two points $P$ and $Q$ of a conic meet at a point $T$, and if the straight line $PQ$ meets the directrix corresponding to focus $S$ in a point $K$, then show that the angle $\angle KST$ is a right angle. \pts{6}
\end{parts}

\medskip\rule{\linewidth}{0.2pt}

\medskip


\noindent   \textbf{Question 3.}

\begin{parts}
  \item Find the equation of the plane which passes through the points $(-1, 1, 1)$ and $(1, -1, 1)$ and is perpendicular to the plane $x + 2y + 2z = 5$. \pts{6}
  \item Show that the lines:
  \[
  \frac{x-1}{2} = \frac{y-2}{3} = \frac{z-3}{4} \quad  \text{and} \quad \frac{x-4}{4} = \frac{y-3}{2} = \frac{z-2}{-1}
  \]
  are coplanar, and find the equation of the plane containing them. \pts{6}
\end{parts}

\medskip\rule{\linewidth}{0.2pt}

\medskip


\noindent   \textbf{Question 4.}

\begin{parts}
  \item Find the coordinates of the points on the sphere $x^2 + y^2 + z^2 - 4x + 2y = 4$ where the tangent planes are parallel to the plane $2x - y + 2z = 1$. \pts{6}
  \item Find the equation of the sphere which passes through the two points $(0, 3, 0)$, $(-2, -1, -4)$ and cuts orthogonally the two spheres:
  \[
  x^2 + y^2 + z^2 - 2x + 3z - 2 = 0 \quad  \text{and} \quad 2(x^2 + y^2 + z^2) + x + 3y + 4 = 0
  \] \pts{6}
\end{parts}

\medskip\rule{\linewidth}{0.2pt}

\medskip


\noindent   \textbf{Question 5.}

\begin{parts}
  \item Prove that the equation $a x^2 + b y^2 + c z^2 + 2ux + 2vy + 2wz + d = 0$ represents a cone if:
  \[
  \frac{u^2}{a} + \frac{v^2}{b} + \frac{w^2}{c} = d
  \] \pts{6}
  \item Find the condition that the plane $ux + vy + wz = 0$ may touch the cone $a x^2 + b y^2 + c z^2 = 0$. \pts{6}
\end{parts}

\medskip\rule{\linewidth}{0.2pt}

\medskip


\noindent   \textbf{Question 6.}

\begin{parts}
  \item Find the equation of the cylinder whose generators are parallel to the line $\frac{x}{1} = \frac{y}{2} = \frac{z}{3}$ and whose guiding curve is the ellipse $x^2 + 2y^2 = 1, z = 0$. \pts{6}
  \item Find the equation of a right circular cylinder of radius $2$ whose axis is the line:
  \[
  \frac{x-1}{2} = \frac{y-2}{3} = \frac{z-3}{1}
  \] \pts{6}
\end{parts}

\medskip\rule{\linewidth}{0.2pt}

\medskip


\noindent   \textbf{Question 7.}

\begin{parts}
  \item A point $P$ moves so that the section of the enveloping cone of the ellipsoid $\frac{x^2}{a^2} + \frac{y^2}{b^2} + \frac{z^2}{c^2} = 1$, with $P$ as vertex, by the plane $z = 0$ is a circle. Show that the locus of $P$ is:
  \[
  \frac{y^2(b^2-a^2)}{b^4} + \frac{z^2(c^2-a^2)}{c^4} = 1, \ x = 0 \quad  \text{and} \quad \frac{x^2(a^2-b^2)}{a^4} + \frac{z^2(c^2-b^2)}{c^4} = 1, \ y = 0
  \] \pts{6}
  \item Show that the plane $8x - 6y - z = 5$ touches the paraboloid $\frac{x^2}{2} - \frac{y^2}{3} = z$ and find the coordinates of the point of contact. \pts{6}
\end{parts}

\vfill
\rule{\linewidth}{0.4pt}

\begin{center}\small
  \href{https://bhu-exam.vercel.app/}{Click here to visit our website}\\[4pt]
  \href{https://me-aryan.vercel.app/}{Click here to visit our contributor}\\[4pt]
  \href{https://quashan-me.vercel.app/}{Click here to visit our contributor}
\end{center}

\end{document}
